\section{Результаты прямой и обратной задач}
\label{sec:results}

\subsection{Прямая задача}
\draft{Для data1 доступно опубликованное эталонное решение. Сравнение
выполнено на исходной геометрии XLSX по значениям на кровле, подошве
и средним по слою. На рис.~\ref{fig:data1-temperature}
и~\ref{fig:data1-ro} показаны современные поля и знаковые разности.
Используется относительная табличная норма без пространственных весов:}
\begin{equation}
 \varepsilon_{\mathrm{tab}}(u)=100
 \frac{\left(\sum_j(u_j^{\mathrm{calc}}-u_j^{\mathrm{ref}})^2\right)^{1/2}}
      {\left(\sum_j(u_j^{\mathrm{ref}})^2\right)^{1/2}}\quad[\%].
\end{equation}
\reportfigure{12_data1_temperature.png}
  {Data1, современная температура: эталон, расчёт и расчёт минус эталон.
  Поля восстановлены по значениям на кровле и подошве каждого слоя.}
  {fig:data1-temperature}
\reportfigure{13_data1_reflectance.png}
  {Data1, современная зрелость Ro: эталон, расчёт и их знаковая разность.}
  {fig:data1-ro}

\draft{По средним значениям слоёв ошибки T/Ro за все 18 возрастов равны
2,803/1,859\%; для современной конфигурации --- 3,598/2,385\%.
Рисунок~\ref{fig:profiles} показывает вертикальные профили у области
наибольшего расхождения и вдали от неё. На рис.~\ref{fig:data2}
показан второй прямой набор, для которого эталон не опубликован.}
\reportfigure{14_data1_vertical_profiles.png}
  {Data1: вертикальные профили T и Ro при $x=282$ и $1400$ м;
  синие линии --- расчёт, оранжевые --- эталон.}
  {fig:profiles}
\reportfigure{15_data2_fields.png}
  {Data2: современные поля T и Ro с границами слоёв.}
  {fig:data2}
\FloatBarrier

\subsection{Обратная задача}
\draft{Начальный и найденный потоки показаны на рис.~\ref{fig:inverse-q}.
Изменение потока на последней эпохе требует дальнейшей проверки весов
и размерности параметризации; найденная история не объявляется единственной.}
\reportfigure{16_inverse_basal_flux.png}
  {Начальное приближение и найденный нижний поток в 17 узлах
  с линейной интерполяцией между ними.}
  {fig:inverse-q}

\draft{Подбор выполнен на сетке $200\times100$ м: 10 итераций,
44 прямых вычисления, $J=265{,}883$; остановка по лимиту итераций.
Рисунок~\ref{fig:optimization} показывает лучшее значение функционала
и стоимость начальной серии из 18 вычислений. Контроль найденного потока
на сетке $200\times50$ м даёт $J_T=205{,}898$, $J_{\Ro}=55{,}608$.}
\reportfigure{17_inverse_objective.png}
  {История функционала по числу прямых вычислений; выделена начальная серия оптимизатора.}
  {fig:optimization}

\draft{Итоговые профили по пяти скважинам приведены на
рис.~\ref{fig:wells-temperature} и~\ref{fig:wells-ro}.
В допуске 456/500 температур и 39/50 значений Ro.
Среднеквадратичные относительные невязки --- 1,283\% и 2,109\%,
максимальные --- 6,743\% и 6,907\% соответственно.}
\reportfigure{18_inverse_wells_temperature.png}
  {Температура в скважинах A--E при найденном потоке:
  расчёт, измерения и интервалы погрешности 2\%.}
  {fig:wells-temperature}
\reportfigure{19_inverse_wells_reflectance.png}
  {Зрелость Ro в тех же скважинах: расчёт, измерения и интервалы погрешности 2\%.}
  {fig:wells-ro}

\draft{Рисунок~\ref{fig:inverse-fields} показывает современные T и Ro
во всём разрезе при том же потоке, включая участки вне скважин.
Для этих полей эталон неизвестен.}
\reportfigure{20_inverse_fields.png}
  {Полные современные поля T и Ro обратного набора при найденном потоке.}
  {fig:inverse-fields}
